\addvspace{1cm}
\subsection{Eigenschaften der einseitigen z-Transformation}

\textbf{Allgemeine Eigenschaften}
Die Transformierte hat zunächst die Form einer Potenzreihe der komplexen Variablen $\frac{1}{z}$, wobei die Koeffizienten genau die Abtastwerte $x[k]$ sind. 
Aus den Eigenschaften der Potenzreihe ergibt sich, dass es einen Kovergenzradius $r$ gibt, sodass $X(z)$ für $|z|>\frac{1}{r}$ konvergiert und für $|z|<\frac{1}{r}$ divergiert \cite[S.233]{book2}.

\addvspace{5cm}
\textbf{Abbildung der s-Ebene auf die z-Ebene}

Da die z-Transformierte eines diskreten Signals direkt aus der Laplace-Transformierten des zeitkontinuierlichen Signals ableitet, ist es sinnvoll, zu betrachten, wie sich Punkte wie Pol- und Nullstellen aus der s-Ebene auf die z-Ebene übertragen.

Die Transformationsgleichungen von der s-Ebene in die z-Ebene ergeben sich wie folgt:

\begin{align}
    \notag
    s = \sigma + j \omega \text{ und }& \\
    \notag
    z = \xi + j \eta  = e^{sT}& = e^{\sigma + j \omega} = e^{\sigma} \cdot e^{j \omega} = r \cdot e^{j \omega}\\
     \notag
    \xi + j \eta = e^{(\sigma + j \omega)T} & = e^{\sigma T} \cdot [\cos(\omega T) + j \sin(\omega T)] \\
     \notag
    \xi = e^{\sigma T} \cdot \cos(\omega T) &\text{ und } \eta = e^{\sigma T} \cdot \sin(\omega T) \\
     \notag
\end{align}

Der Radius $r$ beschreibt dabei den Kreis um den Ursprung auf dem die z-Transformierte mit konstantem Realteil von $z$ liegt.

\begin{itemize}
    \item Punkte mit konstantem Realteil $\sigma$ liegen auf einem Kreis mit dem Radius $\rho=e^{\sigma T}$.
    \item die $j \omega$-Achse der s-Ebene wird $2\pi$-periodisch auf den Einheitskreis der z-Ebene aufgewickelt.
    \item Die Stellen $s = 0,\pm \frac{j2\pi}{T},\pm \frac{j4\pi}{T},...$ werden auf den Punkt $z=1$ abgebildet.
    \item Die Stellen $s = \pm \frac{j\pi}{T},\pm \frac{j3\pi}{T},...$ werden auf den Punkt $z=-1$ abgebildet.
    \item die rechte s-HE mitsamt Polen und Nullstellen wird auf den äußeren Bereich des Einheitskreises abgebildet.
    \item die linke s-HE mitsamt Polen und Nullstellen wird auf den inneren Bereich des Einheitskreises abgebildet.
\end{itemize}
\cite[S.234]{book2}

%((Bsp-Bild aus Matlab einfügen))

\addvspace{1cm}
\textbf{Region of convergence (ROC)}

Damit die Transformierte konvergiert, muss die Summe der exponentiell gewichteten Abtastwerte absolut summierbar sein. Es muss also gelten:

\begin{equation}
    \sum_{k=0}^{\infty} |x[k]| \cdot |z|^{-k} = \sum_{k=0}^{\infty} |x[k] \cdot r^{-k}| < \infty
    \label{eq:absKonvergenz}
\end{equation}

Da die DFT den Faktor $r^{-k}$ nicht enthält, konvergiert sie nur, wenn die Abtastwerte an sich absolut summierbar sind. Im Gegensatz dazu konvergiert die z-Transformation mit passendem $r$ auch für Abtastwerte von Funktionen, die langsamer als die Exponentialfunktion wachsen \cite[S.182]{book3}.

Die ROC ist eine Kreisscheibe, die alle Radien $|z|$ umfasst, für die die oben genannte Summe konvergiert. Je höher der Radius ist, desto größer muss das abgetastete Signal gedämpft werden. An dieser Eigenschaft lässt sich auch erkennen, dass Systeme mit einer ROC, die außerhalb des Einheitskreises liegt, instabil sind.
In manchen Fällen, kann die ROC auch bis $|z| \rightarrow  \infty$ oder bis zum Ursprung $|z|=0$ ausgedehnt sein.

Wichtig zu erwähnen ist, dass die ROC keine Polstellen enthalten kann und immer ein zusammenhängender Bereich ist. Dabei wird die ROC durch die Polstellen begrenzt.

Für die z-Transformierte einer rechtsseitigen Folge, bei der $x[k] = 0$ für $k<0$ und die für $n \leq N < \infty$ nicht verschwindet, erstreckt sich die ROC von der äußersten Polstelle bis zu $\infty$.
(Bei endlichen Folgen ist die ROC der zugehörigen z-Transformierten immer die gesamte z-Ebene mit möglichen Ausnahmen bei $z=0$ und $z=\infty$.) \cite[S.191]{book3}

